\usepackage{amsmath,amssymb,stmaryrd,xspace}
\usepackage{amsthm}
\usepackage{tikz-cd}
\usepackage{multirow}
\usepackage{mathtools}
\usepackage{url}
% \usepackage{makeidx}
\usepackage{relsize}
\usepackage{bm} %For bold math.
% \usepackage[inline]{enumitem}
% \setenumerate{label=\arabic*)}
% \setenumerate{label=\arabic*.}
\newcounter{counter}
% \input{./style-options}

\newcommand{\nonext}[1]{#1'}

\newcommand{\tikzfig}[1]{\scalebox{1.0}{\input{./figures/#1.tikz}}} %was 0.8
\newcommand{\tikzitcd}[1]{\scalebox{1.0}{\input{./figures/#1.tikz}}} 
\newcommand{\boxpic}[1]{\scalebox{0.6}{\input{./figures/#1.tikz}}}
\newcommand{\mcatfig}[1]{\scalebox{0.6}{\input{./figures/#1.tikz}}}
\newcommand{\storecounter}{\setcounter{counter}{\value{enumi}}}
\newcommand{\resumecounter}{\setcounter{enumi}{\value{counter}}}


\newcommand{\DetSt}{\St}
\newcommand{\ProbSt}{\Prob}
\newcommand{\phimax}{\varphi^{\textrm{max}}}

% \newcommand{\bibxiv}[1]{\url{#1}} % Wrapper around arxiv addresses in note field of bibtex
\newcommand{\bibxiv}[1]{{#1}} % Wrapper around arxiv addresses in note field of bibtex

\newcommand{\indef}[1]{\deff{#1}} %in-line definition
\newcommand{\emps}[1]{\emph{#1}} %emphasis
\newcommand{\deff}[1]{{\boldmath\textbf{#1}}} %definition inside a def
\newcommand{\Expcat}{\cat{Exp}}
\newcommand{\Pexpcat}{\cat{PExp}}
\newcommand{\PExp}{\mathbb{PE}}
\newcommand{\Exp}{\mathbb{E}}
\newcommand{\WExp}{\mathbb{\bar{E}}}
\newcommand{\Pexpi}{\overline{\Pexpcat}}
\newcommand{\Sys}{\cat{Sys}}
\newcommand{\rep}{\mathsf{r}} %Generic symbol for a repertoire
% \newcommand{\Exp}{\mathsf{IIT}}

\newcommand{\BasicIIT}{\mathsf{IIT}}
\newcommand{\ClassIIT}{\BasicIIT_{\mathsf{3.0}}}

%\newcommand{\IIT}{\mathsf{IIT}}
\newcommand{\QShape}{\mathbb{Q}}
\newcommand{\cut}[2]{{#1}^{#2}}

\newcommand{\temp}[1]{\color{red}{#1}}

\newcommand{\marg}[2]{{#1}|_{#2}} %marginal of a state


% \newcommand{\decomppic}[3]{\tikzfig{decomp-macro}}

\newenvironment{exampleslist}
    {\begin{enumerate}[leftmargin=*,label=\arabic*., ref=\arabic*]}
    {\end{enumerate}}

% \newenvironment{exampleslistresume}
%     {\begin{enumerate}[leftmargin=*,label=\arabic*.,resume]}
%     {\end{enumerate}}

\newenvironment{theoremlist}
    {\begin{enumerate}[leftmargin=*,label=\roman*]}
    {\end{enumerate}}

\newenvironment{inlinelist}
   {\begin{enumerate*}[label=\arabic*), ref=\arabic*)]}
    {\end{enumerate*}}


%Klenne equality
\newcommand{\kleene}{\simeq}

% % Indexing indecision in an operational theory. Allows for swapping to I, J notation if desired.
% \newcommand{\x}{i}
% \newcommand{\y}{j}
% \newcommand{\z}{k}
% \newcommand{\w}{l}
% \newcommand{\oi}{j} %SET TO j if \x is i, i if \x is x
% \newcommand{\X}{X}
% \newcommand{\Y}{Y}


% \newcommand{\effpullbone}{\text{(a)}}
% \newcommand{\effpullbtwo}{\text{(b)}}


% %Environments
% \theoremstyle{plain}
% %\newtheorem{assumption}{Assumption}[section]
% \newtheorem{assumption}{Axiom}
% \newtheorem{principle}{Principle}
% \newtheorem{property}{Property}

% \newtheorem{theorem}{Theorem}[chapter]
% \newtheorem{proposition}[theorem]{Proposition}
% \newtheorem{corollary}[theorem]{Corollary}
% \newtheorem{lemma}[theorem]{Lemma}
% \newtheorem{conjecture}[theorem]{Conjecture}
% \theoremstyle{definition}
% \newtheorem{definition}[theorem]{Definition}
% \newtheorem{example}[theorem]{Example}
% \newtheorem{examples}[theorem]{Examples}
% \newtheorem{remark}[theorem]{Remark}
% \newtheorem{question}[theorem]{Question}



%Including images, notes to self 
\newcommand{\sketch}[1]{\includegraphics[scale=0.3]{#1}}
\newcommand{\notetoself}[1]{}
\newcommand{\omitfornow}[1]{}
\newcommand{\omitthis}[1]{}
\newcommand{\reconstructionsection}[1]{#1} %Should leave out of here eventually
\newcommand{\phasedcoprodsection}[1]{#1}%{#1}
\newcommand{\ignore}[1]{#1}
\newcommand{\TODO}[1]{\marginpar{\footnotesize \textbf{TODO:} #1}}

\newcommand{\todo}[1]{\marginpar{\footnotesize #1}}

%Categories
\newcommand{\Set}{\cat{Set}}
\newcommand{\VecSp}{\cat{Vec}}
\newcommand{\VecC}{\VecSp_\mathbb{C}}
\newcommand{\Veck}{\VecSp_\fieldk}
% \newcommand{\FVecSp}{\cat{FVec}}
\newcommand{\FVeck}{\cat{FVec}_\fieldk}
\newcommand{\FHilb}{\cat{FHilb}}
\newcommand{\Hilb}{\cat{Hilb}}
\newcommand{\Rel}{\cat{Rel}}
\newcommand{\Prob}{\mathcal P}
\newcommand{\Grp}{\cat{Grp}}
\newcommand{\PFun}{\cat{PFun}}
\newcommand{\KlDT}{\cat{Class}}
\newcommand{\KlSD}{\cat{Class_{\mathsf{p}}}}
\newcommand{\KlD}{\Kleisli{D}}
\newcommand{\KlMn}{\Kleisli{\multiset_\mathbb{N}}}
\newcommand{\Class}{\cat{Class}}
\newcommand{\Classm}{\cat{Class}_\mathsf{m}} % Variation of class in which objects are finite metric spaces.
\newcommand{\Kleisli}[1]{\mathsf{Kl}(#1)}
\newcommand{\Stoch}{\cat{Stoch}}
\newcommand{\CStarop}{\cat{CStar}^\op}
\newcommand{\CStarSUop}{\cat{CStar}^\op_{\mathsf{su}}}
\newcommand{\CStarUop}{\cat{CStar}^\op_{\mathsf{u}}}
\newcommand{\FCStar}{\cat{CStar}}
\newcommand{\vNop}{\cat{vNA}^\op}
\newcommand{\vNSUop}{\cat{vNA}^\op_{\mathsf{su}}}
\newcommand{\vNUop}{\cat{vNA}^\op_{\mathsf{u}}}
\newcommand{\MultSet}{\cat{MultSet}}
\newcommand{\Quant}[1]{\cat{Quant}_{#1}}
\newcommand{\QuantSU}{\cat{Quant}_{\mathsf{sub}}}
\newcommand{\HilbP}{\Hilb_\sim}
\newcommand{\FHilbP}{\FHilb_\sim}
\newcommand{\Mat}{\cat{Mat}} %NOte this was in mathsf font for a long time. 
\newcommand{\MatS}{\Mat_S}
\newcommand{\Rep}{\cat{Rep}} %Category of representations of a group in FHilb.
\newcommand{\Spek}{\cat{Spek}}
\newcommand{\MSpek}{\cat{MSpek}}
\newcommand{\SpekOrMSpek}{\cat{(M)Spek}}
\newcommand{\IV}{IV}

% \newcommand{\MatR}{\Mat_R}
\newcommand{\PhaseMat}[2]{\cat{Mat}_{#1}^{#2}}
\newcommand{\GSet}{\cat{G}\text{-}\Set}
\newcommand{\VecP}{{\VecSp{}}_{\sim}}
\newcommand{\FVecP}{\cat{FVec}_{\sim}}
\newcommand{\VecProj}{\cat{Proj}_k} %\cat{Vec}_{\mathsf{Proj}}}
% 
% \newcommand{\DComp}{\text{DComp}}
\newcommand{\Stab}{\cat{Stab}}
\newcommand{\Pos}{\cat{Pos}}
\newcommand{\PreOrd}{\cat{PreOrd}}
\newcommand{\RelC}{\cat{Rel}(\catC)}
\newcommand{\RelCAlt}{\cat{Rel}(\catC)^{\bot}}
\newcommand{\RelAlt}{\cat{Rel}^{\bot}}
\newcommand{\catG}{\cat{G}} 
\newcommand{\ParTestCat}{\cat{ParTest}}
\newcommand{\boolB}{\mathbb{B}}

\newcommand{\CMon}{\cat{CMon}} 
\newcommand{\DCM}{\cat{DCM}} 
\newcommand{\PCM}{\cat{PCM}}
\newcommand{\PCMD}{\cat{Par}} %Todo: should somehow add discard symbols here. 
\newcommand{\CMonD}{\cat{Tot}}

\newcommand{\Concepts}{\mathbb{C}}

% \newcommand{\EMod}{\cat{EMod}}
% \newcommand{\CStarPU}{\cat{CStar}_{PU}}
% \newcommand{\CStarCPU}{\cat{CStar}_{\mathsf{CPU}}}
% \newcommand{\OAb}{\cat{OAb}}
% \newcommand{\OUG}{\cat{OUG}}
% \newcommand{\OUS}{\cat{OUS}}
% \newcommand{\CStarMIU}{\cat{CStar}_{MIU}}

% \newcommand{\CP}{\ensuremath{\cat{CP}\xspace}}
% \newcommand{\CPs}{\ensuremath{\CP^*}\xspace}
\newcommand{\CPM}{\ensuremath{\mathsf{CPM}}\xspace}
\newcommand{\Split}[1]{\ensuremath{\mathrm{Split}}\xspace}
\newcommand{\Splitcaus}{\ensuremath{\mathrm{Split}_{c}}\xspace}
% \newcommand{\Splitd}{\ensuremath{\mathrm{Split}^\dag}\xspace}
% \newcommand{\V}{\cat{V}}
% \newcommand{\C}{\ensuremath{\mathbb{C}}}
% \newcommand{\CCStar}{\cat{CCStar}}
% \newcommand{\CStarCP}{\cat{CStar}_{CP}}
% \newcommand{\Conv}{\cat{Conv}}
\newcommand{\cat}[1]{\ensuremath{\mathbf{#1}}\xspace}
\newcommand{\catJ}{\cat{J}}
\newcommand{\catD}{\cat{D}}
\newcommand{\catA}{\cat{A}}
\newcommand{\catC}{\cat{C}}
\newcommand{\catB}{\cat{B}}

% \newcommand{\catX}{\cat{X}}
% \newcommand{\catH}{\cat{H}}
% \newcommand{\catF}{\cat{F}}
% \newcommand{\catE}{\cat{E}}
% \newcommand{\catP}{\mathbb{P}}



%Operational Theories
\newcommand{\ClassDet}{\mathsf{ClassDet}}
\newcommand{\ClassProb}{\mathsf{ClassProb}}
\newcommand{\ClassProbCat}{\cat{ClassProb}} %Category of events of ClassProb.
%\newcommand{\Quant}{\mathsf{QuantOp}}
\newcommand{\MatOT}{\mathsf{Mat}} 
% \newcommand{\FinHilbOTC}{\mathsf{Quant}} 
\newcommand{\QuantOT}{\mathsf{Quant}}
\newcommand{\QuantClassOT}{\mathsf{QuantClass}} 
\newcommand{\FinCStarOTC}{\mathsf{FCStar}} 
\newcommand{\CStarOTC}{\mathsf{CStar}} 
\newcommand{\CStarOT}{\mathsf{CStar}} 
\newcommand{\RelOT}{\mathsf{Rel}}
\newcommand{\Kl}{\cat{Kl}} %TODO: SORT OUT THIS KLEISLI NOTATION
\newcommand{\KleisliD}{\cat{Kl}(\mathcal{D})} %Sort out Kleisli stuff 
% \newcommand{\FDimCStar}{\cat{FdimCStar}_{\text{u}}^{\op}}
% \newcommand{\FDimCStarSU}{\cat{FdimCStar}_{\text{su}}^{\op}}
% \newcommand{\FDimCStarBracket}{\cat{(Fdim)CStar}_{\text{u}}^{\op}}
% \newcommand{\FDimCStarSUBracket}{\cat{(Fdim)CStar}_{\text{su}}^{\op}}

%Categories of operational theories
\newcommand{\OTCat}{\cat{OT}}
\newcommand{\OTRep}{\OTCat^+}
\newcommand{\OTCatOrRep}{\OTCat^{(+)}}

\newcommand{\OTRepStrict}{\OTRep_{\mathsf{strict}}}
% \newcommand{\OTRep}{\cat{OT}_{\mathsf{rep}}}
% \newcommand{\OTRepStrict}{\OTRep^{\mathsf{strict}}}
\newcommand{\OpCats}{\cat{OCat}} %Category of operational categories.
\newcommand{\OpCatsComp}{\OpCats_{\mathsf{comp}}} %Category of operational categories.
\newcommand{\TestCats}{\cat{TestCat}} %Category of operational categories.


%\newcommand{\OpCat}{\TestCat} %Category of ops/tsts of an OTC
\newcommand{\PTest}{\cat{PTest}} %Category of partial tests of an operational theory. 
\newcommand{\Test}{\cat{Test}} %Category of partial tests of an operational theory. 
\newcommand{\TestCat}{\cat{Test}}
\newcommand{\Event}{\cat{Event}} %Category of events of an OPT 
%\newcommand{\ParOpCat}{\cat{ParTest}} %Category of partial ops/tsts of an OTC
%\newcommand{\OpStructure}{\mathsf{Test}} %Op/tst structure of an OTc





%Subcategories
\newcommand{\pos}{\mathrm{pos}} %Can do mathrm or mathsf 
\newcommand{\causal}{\mathsf{caus}}
%\newcommand{\supercausal}{\mathsf{sc}}
\newcommand{\subcausal}{\mathsf{sc}}%{\mathsf{subcausal}}
\newcommand{\pure}{\mathsf{pure}}
\newcommand{\prepure}{\mathsf{p}} %Notation for pure morphisms of general environment structure.
\newcommand{\sa}{\mathsf{sa}} %Self adjoint
\newcommand{\Rplus}{\mathbb{R}^+} %Positive real numbers 
\newcommand{\Rpos}{\Rplus} %Positive real numbers 
\newcommand{\Qpos}{\mathbb{Q}^+} %Positive rationals numbers 


\newcommand{\dkernel}{kernel} %Notion of kernel purity in dagger compact category. Either kernel or \dagger-kernel

%Category theory
\newcommand{\id}[1]{\ensuremath{\mathrm{id}_{#1}}}
\newcommand{\Sub}{\mathrm{Sub}}
\newcommand{\Cut}{\mathrm{Cut}}
\newcommand{\Aut}{\mathrm{Aut}} %Automorphisms
\newcommand{\ob}[1]{\ensuremath{\mathrm{ob}({#1})}}
\newcommand{\obC}{\ob{\catC}}
\newcommand{\op}{\ensuremath{\mathrm{\rm op}}}
\DeclareMathOperator{\Ob}{Ob}
\DeclareMathOperator{\Mor}{Mor}
\DeclareMathOperator{\dom}{dom}
\DeclareMathOperator{\cod}{cod}
\newcommand{\Multicat}{\mathsf{M}} %Multicategory of a monoidal category
\newcommand{\mcat}[1]{\cat{#1}} %Font for Multicategories 
\newcommand{\partialcomp}{\hspace{3.5pt} \partialsymbolcomp \hspace{2.5pt}} %Composition in Kleisli
\newcommand{\tikzcdpartialright}[1]{\arrow{r}[description, inner sep = 0]{\partialsymbollargearrow}[above]{{#1}}}



\newcommand{\partialsmallarrow}[1][1.2pt]{ 
\begin{tikzpicture}[baseline=-0.58ex]
\draw[black, fill=white,radius={#1}] (0,0) circle ; 
\draw[black, fill=white,radius=0.1pt] (0,0) circle ; 
\end{tikzpicture}
}
\newcommand{\circlearrow}[1]{\ensuremath{\mathrlap{\hspace{2.5pt}\to}{\hspace{6pt}{#1}\hspace{6pt}}}}

\newcommand{\partialto}{\circlearrow{\partialsmallarrow[2.5pt]}}

\newcommand{\partialsymbolcomp}{ 
\begin{tikzpicture}[baseline=-0.5ex]
\draw[black, fill=white,radius=3.6pt] (0,0) circle ; 
\draw[black, fill=white,radius=0.6pt] (0,0) circle ; 
\end{tikzpicture}
}
\newcommand{\partialsymbollargearrow}{ 
\begin{tikzpicture}[baseline=-0.6pt]
\draw[black, fill=white,radius=3pt] (0,0) circle ; 
\draw[black, fill=white,radius=0.2pt] (0,0) circle ; 
\end{tikzpicture}
}


%Discarding and tests
%NEW DISCARD SYMBOL

\newcommand{\discard}[1]{\ensuremath{\tinygroundnew_{#1}}}
\newcommand{\discardold}[1]{\ensuremath{\tinyground_{#1}}}
\newcommand{\discardflip}[1]{\ensuremath{\tinygroundflipnew_{#1}}}
\newcommand{\maxmix}[1]{\ensuremath{\tinygroundflipnewer_{#1}}}
% \newcommand{\uniform}[1]{\ensuremath{\uniformprob_{#1}}}
\newcommand{\uniform}[1]{{\omega_{#1}}} %Sean: I think this is simpler
%\newcommand{\discardflip}[1]{{\discard{#1}}^{\dagger}}
%\newcommand{\bigovee}{\mathlarger{\ovee}}
\newcommand{\ts}[1]{(#1)} %Test/partial test notation

\DeclareFontFamily{U}{mathux}{\hyphenchar\font45}
\DeclareFontShape{U}{mathux}{m}{n}{
      <5> <6> <7> <8> <9> <10>
      <10.95> <12> <14.4> <17.28> <20.74> <24.88>
      mathux10
      }{}
\DeclareSymbolFont{mathux}{U}{mathux}{m}{n}

\DeclareMathSymbol{\bigovee}{1}{mathux}{"8F}

\DeclareMathSymbol{\bigperp}{1}{mathux}{"4E}


% This is all to define the special `total as partial' brackets
\DeclareFontFamily{U}{mathb}{\hyphenchar\font45}
\DeclareFontShape{U}{mathb}{m}{n}{
      <5> <6> <7> <8> <9> <10> gen * mathb
      <10.95> mathb10 <12> <14.4> <17.28> <20.74> <24.88> mathb12
      }{}
\DeclareSymbolFont{mathb}{U}{mathb}{m}{n}
\DeclareFontSubstitution{U}{mathb}{m}{n}
\DeclareMathSymbol{\mylgroup}{\mathbin}{mathb}{'160}
\DeclareMathSymbol{\myrgroup}{\mathbin}{mathb}{'161}

\newcommand{\totaspar}[1]{ \mylgroup {#1} \myrgroup } %Viewing a total arrow as a partial one 




%Math and Math Operators
\newcommand{\eqdef}{\vcentcolon=} %Use instead of :=
\newcommand{\eqdefrev}{=\vcentcolon}
\newcommand{\curly}[1]{\ensuremath{\mathcal{#1}}}
% \newcommand{\supp}{\mathsf{supp}}
\newcommand{\dring}[1]{D(#1)}%{\hat{#1}} %Notation for difference ring
% \newcommand{\norm}[1]{||#1||}
\newcommand{\infinity}{\infty}
\newcommand{\Sup}{\bigvee}
\newcommand{\curlyS}{\curly{S}}
\newcommand{\curlyI}{\curly{I}}
\newcommand{\curlyT}{\curly{T}}
\newcommand{\sem}[1]{\ensuremath{\llbracket #1 \rrbracket}}
\newcommand{\defined}{\ensuremath{\mathop{\downarrow}}}
\newcommand{\pfun}{\relto} %Partial function arrow
\newcommand{\curlyB}{B} %For bounded operators on a Hilbert space. 
\newcommand{\fieldk}{k} %A field called k

%Quantum
\newcommand{\inprod}[2]{\ensuremath{\langle #1\,|\,#2 \rangle}}
% \DeclareMathOperator{\Tr}{Tr}
%\newcommand{\Dbl}{\mathsf{Dbl}} %Doubling functor C -> CPM(C)
\newcommand{\Dbl}[1]{\widehat{#1}}
%{\mathsf{Dbl}(#1)} %Doubling functor C -> CPM(C)
% \newcommand{\bra}[1]{\ensuremath{\langle #1 |}}
% \newcommand{\ket}[1]{\ensuremath{| #1 \rangle}}


%Effectus
\newcommand{\asrt}{\text{asrt}}
\newcommand{\Pred}{\text{Pred}}
\newcommand{\Eff}{\text{Eff}}
\newcommand{\Stat}{\text{Stat}}

%Arrows
\newcommand{\isomto}{\xrightarrow{\sim}}


%Totalisation
\newcommand{\To}{\mathbb{T}}
\newcommand{\Ro}{R} %Alternative form of totalisation
\newcommand{\multiset}{\mathcal{M}}


%Dagger Kernels
\newcommand{\Ker}{\mathrm{Ker}}
\newcommand{\Coker}{\mathrm{Coker}}
\newcommand{\kernelto}{\rightarrowtail}
\newcommand{\img}{\mathrm{im}}
\newcommand{\Img}{\mathrm{Im}}
\newcommand{\CoIm}{\mathrm{Coim}}
\newcommand{\coim}{\mathrm{coim}}
\newcommand{\coker}{\mathrm{coker}}
\newcommand{\DKer}{\mathrm{DKer}} %Collection of dagger kernels. 

%Encoding and Decoding maps and object of ideal compression
\newcommand{\enc}{E}
\newcommand{\dec}{D}
\newcommand{\idob}{F}


%Extension maps (purification)
%Replacing this notation with `minimal dilations'. 
%\newcommand{\ext}[1]{\mathrm{ext}(#1)}
%\newcommand{\Ext}[2]{{#1}_{#2}} %\ext{A}{e} for e : A -> I
\newcommand{\ext}[1]{\mdil{#1}}
\newcommand{\Ext}[2]{\MDil{#2}}%{#1}_{#2}} %\ext{A}{e} for e : A -> I
\newcommand{\MDil}[1]{M({#1})} %for f : A -> B
\newcommand{\mdil}[1]{\mathrm{min}(#1)}
\newcommand{\extensions}{extensions}
\newcommand{\extension}{extension}
\newcommand{\med}{\nu} %Mediating map between coimages and min dilations. 


%Limits and Phased Limits
\newcommand{\coproj}{\kappa}
\newcommand{\biprod}{\oplus}
\newcommand{\pinit}{0}%{\dot{0}} %Notation for a phased initial object.
\newcommand{\pterm}{\dot{1}} %Notation for a phased terminal object.
\newcommand{\pcoprod}{\mathbin{\dot{+}}}
%\newcommand{\pcoprod}{\widetilde{+}}%{+^p}
%\newcommand{\pcoprod}{\boxplus}%{+^p}
%\newcommand{\pcoprod}{\widetilde{+} }%{+^p}
\newcommand{\pprod}{\mathbin{\dot{\times}}}
%\newcommand{\pbiprod}{\widetilde{\biprod}}
\newcommand{\pbiprod}{\mathbin{\dot{\oplus}}}
\newcommand{\pcprod}{\mathbin{\dot{\times}}}
\newcommand{\pcoproj}{\coproj}
\newcommand{\pproj}{\pi}
\newcommand{\simP}{\sim} %equivalence relation of monoidal category `up to global 'phases P
\newcommand{\quotP}{\mathbb{P}} %\catC_\quotP is quotient of \catC by \simP
\newcommand{\quot}[2]{{#1} \!/\! {#2}}
\newcommand{\Gen}{I} %General symbol used for generator object in context of phased coproducts.
\newcommand{\IsomGen}{I} %General symbol used for generator object in context of phased coproducts.
\newcommand{\TrI}{\mathbb{T}} %Notation for trivial isomorphisms on object A
\newcommand{\classP}{\mathbb{P}} %equivalence relation of monoidal category `up to global 'phases P


%Phased Limit Constructions
\newcommand{\GPa}[1]{\mathrm{GP}({#1})} %alternative version of GP.
\newcommand{\plusI}[1]{\mathsf{GP}({#1})}
\newcommand{\plusIdag}[1]{\mathsf{GP}^{\dagger}({#1})}
\newcommand{\phase}[1]{\plusI{#1}} %Alternative notation for plusI
\newcommand{\plusconstr}[1]{{#1} \pcoprod {#1}}
\newcommand{\tens}{\hat{\otimes}}
\newcommand{\dtens}{\tens} %Version of $\tens$ used in TikZit diagrams. 
\newcommand{\obb}[1]{{\bm{#1}}}
% \newcommand{\obb}[1]{{\underline{#1}}}
\newcommand{\morr}[1]{\obb{#1}}
\newcommand{\aalpha}{\obb{\alpha}}
\newcommand{\llambda}{\obb{\lambda}}
\newcommand{\ssigma}{\obb{\sigma}}
\newcommand{\rrho}{\obb{\rho}}


% \newcommand{\morr}[1]{\mathbf{#1}}
% \newcommand{\aalpha}{\bm{\alpha}}
% \newcommand{\llambda}{\bm{\lambda}}
% \newcommand{\ssigma}{\bm{\sigma}}
% \newcommand{\rrho}{\bm{\rho}}

%Categories for Universal version of Phased Limits
\newcommand{\MD}{\cat{Co}} %Finite monic distributive coproducts and strict monoidal functors.
\newcommand{\BMD}{\cat{BCo}} %Finite braided monic distributive coproducts and strict monoidal functors.
\newcommand{\MDG}{\cat{Co}_{\mathsf{G}}} %Finite monic distributive coproducts and strict monoidal functors with choise of global phase group.
\newcommand{\BMDG}{\cat{BCo}_{\mathsf{G}}} %Finite monic distributive coproducts and strict monoidal functors with choise of global phase group.
\newcommand{\MDPh}{\cat{PhCo}} %Finite monic distributive phased coproducts and strict monoidal functors.
\newcommand{\BMDPh}{\cat{BPhCo}} %Finite monic distributive phased coproducts and strict monoidal functors.




%Operational Categories
\newcommand{\ptest}{\bot}
\newcommand{\Stc}{\St_\mathsf{c}}
\newcommand{\St}{\text{St}}
\newcommand{\StR}{\text{St}_R}
\newcommand{\EffR}{\text{Eff}_R}
\newcommand{\OTCs}{\cat{(Mon)OTC}}
\newcommand{\MonOTCs}{\OTC}
\newcommand{\OTCPlus}{\cat{OTC^{\otcpllus}}}
\newcommand{\OTCPlusStrict}{\cat{(Mon)OTC^{\otcplus}_{\text{strict}}}}
\newcommand{\Partial}{\mathsf{Par}}
\newcommand{\lift}[1]{ \Partial({ {#1} }) } 
\newcommand{\total}{\text{total}} 
\newcommand{\Events}{\cat{Event}} %Event category of an OTc
\newcommand{\DetEvents}{\cat{DetEvent}} %Event category of an OTc
\newcommand{\opequiv}{\simeq}
\newcommand{\monopequiv}{\simeq_{\otimes}}
\newcommand{\monoropequiv}{\simeq_{(\otimes)}}
\newcommand{\otcplus}{+} %Symbol used for the OTC^+ construction
\newcommand{\OT}{\mathsf{OT_B}} 
%The OT defined from an operational category in partial form
\newcommand{\OTC}{\mathsf{OT}} 
%OTC from strong op cat. 



%Hilbert Spaces
\newcommand{\hilbH}{\mathcal{H}} %Hilbert space
\newcommand{\hilbK}{\mathcal{K}} %Hilbert space
\newcommand{\hilbL}{\mathcal{L}} %Hilbert space
\newcommand{\hilbP}{\mathcal{P}} %Hilbert space

%Ceiling
\DeclarePairedDelimiter{\ceil}{\lceil}{\recil}

%Superpositions Chapter extras
%\newcommand{\tc}{\mathbb{T}} %TC-Eq classes symbol for [-]_\tc 
\newcommand{\tc}{\sim} %TC-Eq classes symbol for [-]_\tc 
\newcommand{\tp}{\mathbb{P}} %Global phase-Eq classes symbol for [-]_\tp
% Should make consistent throughout. Can make either C and D or A and B... 
\newcommand{\cta}{\catA} %Notation for a category which we are about to quotient. 
\newcommand{\ctb}{\catB} %Notation for a category which we see as a quotient. 

%Misc
\newcommand{\evec}{f} %Eve's Channel.

%====================================================
%Diagrams and small symbols, not for use with Tikzit
%====================================================
\usepackage{tikz,xypic}

\usetikzlibrary{decorations.pathreplacing,decorations.markings,arrows.meta,backgrounds}
\pgfdeclarelayer{edgelayer}
\pgfdeclarelayer{nodelayer}
\pgfsetlayers{background,edgelayer,nodelayer,main}


%Careful - these dots are different to those for use with Tikzit. 
\tikzstyle{cdot}=[circle, draw=black, fill=black!25, inner sep=.4ex] %Chris dot style 
\tikzstyle{bigdot}=[dot, inner sep=0pt]
\tikzstyle{whitedot}=[circle, draw=black, fill=white, inner sep=.4ex]
\tikzstyle{greydot}=[circle, draw=black, fill=black!25, inner sep=.4ex] %Added by Sean
\tikzstyle{blackdot}=[circle, draw=black, fill=black, inner sep=.4ex]
\tikzset{arrow/.style={decoration={
    markings,
    mark=at position #1 with \arrow{>[length=2pt, width=3pt]}},
    postaction=decorate},
    reverse arrow/.style={decoration={
    markings,
    mark=at position #1 with {{\arrow{<[length=2pt, width=3pt]}}}},
    postaction=decorate}
}

\newcommand\relto[1][]{\mathbin{\smash{
\begin{tikzpicture}[baseline={([yshift=-1pt]
current bounding box.south)}]
    \node (A) at (0,0) [inner xsep=0pt, inner ysep=1pt, minimum width=0.15cm] {\ensuremath{\scriptstyle #1}};
    \draw [->, line width=0.4pt, line cap=round]
        ([xshift=-2.5pt] A.south west)
        to ([xshift=3pt] A.south east);
    \draw [line width=.4pt] ([xshift=-1pt,yshift=-2pt]A.south) to ([xshift=1pt,yshift=2pt]A.south);
\end{tikzpicture}}}}

\newcommand\sxto[1]{\mathbin{\smash{
\begin{tikzpicture}[baseline={([yshift=-1pt]
current bounding box.south)}]
    \node (A) at (0,0) [inner xsep=0pt, inner ysep=1pt, minimum width=0.15cm] {\ensuremath{\scriptstyle #1}};
    \draw [->, line width=0.4pt, line cap=round]
        ([xshift=-2.5pt] A.south west)
        to ([xshift=3pt] A.south east);
\end{tikzpicture}}}}

\newcommand{\tinymult}[1][cdot]{
\smash{\raisebox{-2pt}{\hspace{-5pt}\ensuremath{\begin{pic}[scale=0.4,yscale=-1]
    \node (0) at (0,0) {};
    \node[#1, inner sep=1.5pt] (1) at (0,0.55) {};
    \node (2) at (-0.5,1) {};
    \node (3) at (0.5,1) {};
    \draw (0.center) to (1.center);
    \draw (1.center) to [out=left, in=down, out looseness=1.5] (2.center);
    \draw (1.center) to [out=right, in=down, out looseness=1.5] (3.center);
    \node[#1, inner sep=1.5pt] (1) at (0,0.55) {};
\end{pic}
}\hspace{-3pt}}}}


\newcommand{\tinydot}[1][cdot]{
\smash{\raisebox{-2pt}{\hspace{-5pt}\ensuremath{\begin{pic}[scale=0.4,yscale=-1]
    \node[#1, inner sep=1.5pt] (1) at (0,0) {};
\end{pic}
}\hspace{-3pt}}}}


\newcommand{\tinymultflip}[1][cdot]{
\smash{\raisebox{-2pt}{\hspace{-5pt}\ensuremath{\begin{pic}[scale=0.4,yscale=1]
    \node (0) at (0,0) {};
    \node[#1, inner sep=1.5pt] (1) at (0,0.55) {};
    \node (2) at (-0.5,1) {};
    \node (3) at (0.5,1) {};
    \draw (0.center) to (1.center);
    \draw (1.center) to [out=left, in=down, out looseness=1.5] (2.center);
    \draw (1.center) to [out=right, in=down, out looseness=1.5] (3.center);
    \node[#1, inner sep=1.5pt] (1) at (0,0.55) {};
\end{pic}
}\hspace{-3pt}}}}



\newcommand{\tinycomult}[1][cdot]{
\smash{\raisebox{-2pt}{\hspace{-5pt}\ensuremath{\begin{pic}[scale=0.4,yscale=1]
    \node (0) at (0,0) {};
    \node[#1, inner sep=1.5pt] (1) at (0,0.55) {};
    \node (2) at (-0.5,1) {};
    \node (3) at (0.5,1) {};
    \draw (0.center) to (1.center);
    \draw (1.center) to [out=left, in=down, out looseness=1.5] (2.center);
    \draw (1.center) to [out=right, in=down, out looseness=1.5] (3.center);
    \node[#1, inner sep=1.5pt] (1) at (0,0.55) {};
\end{pic}
}\hspace{-3pt}}}}

%SEAN: THIS WAS RAISEBOX 1pt 
\newcommand{\tinyunit}[1][cdot]{
\smash{\raisebox{3pt}{\hspace{-3pt}\ensuremath{\begin{pic}[scale=0.4,yscale=-1]
    \node (0) at (0,0) {};
    \node[#1, inner sep=1.5pt] (1) at (0,0.55) {};
    \draw (0.center) to (1.north);
\end{pic}
}\hspace{-1pt}}}}


\newcommand{\tinyunitdual}[1][cdot]{
\smash{\raisebox{1pt}{\hspace{-3pt}\ensuremath{\begin{pic}[scale=0.4,yscale=1]
    \node (0) at (0,0) {};
    \node[#1, inner sep=1.5pt] (1) at (0,-0.55) {};
    \draw[arrow=.9] (0.center) to (1.north);
\end{pic}
}\hspace{-1pt}}}}


\newcommand{\tinycounit}[1][cdot]{
\smash{\raisebox{-5pt}{\hspace{-3pt}\ensuremath{\begin{pic}[scale=0.4,yscale=1]
    \node (0) at (0,0) {};
    \node[#1, inner sep=1.5pt] (1) at (0,0.55) {};
    \draw (0.center) to (1.north);
        \node[#1, inner sep=1.5pt] (1) at (0,0.55) {};
\end{pic}
}\hspace{-1pt}}}}

\newcommand{\tinyunitflip}[1]{\tinycounit{#1}}

% \newcommand{\tinyunitflip}[1][dot]{
% %\smash{\raisebox{1pt}{\hspace{-3pt}\ensuremath{\begin{pic}[scale=0.4,yscale=1]
%  \smash{\raisebox{-4pt}{\hspace{-3pt}\ensuremath{\begin{pic}[scale=0.4,yscale=1]  %Sean: I modified this line for formatting
%     \node (0) at (0,0) {};
%     \node[#1, inner sep=1.5pt] (1) at (0,0.55) {};
%     \draw (0.center) to (1.north);
% \end{pic}
% }\hspace{-1pt}}}}


%was raisebox{-3pt}
\newcommand{\tinycup}{\smash{\raisebox{2pt}{\hspace{-2pt}\ensuremath{\begin{pic}[scale=0.2]%,string]
   \pgftransformscale{1.5} \draw[arrow=.6, scale = 1] (0,0) to[out=-90,in=-90,looseness=1.5] (1.5,0);
\end{pic}}}}}

\newcommand{\tinycap}{\smash{\raisebox{-3pt}{\hspace{-2pt}\ensuremath{\begin{pic}[scale=0.2, yscale=-1]%,string]
   \pgftransformscale{1.5} \draw[arrow=.6, scale = 1] (0,0) to[out=-90,in=-90,looseness=1.5] (1.5,0);
\end{pic}}}}}

\newcommand{\tinycupswap}{\smash{\raisebox{-3pt}{\hspace{-2pt}\ensuremath{\begin{pic}[scale=0.2]%,string]
   \pgftransformscale{1.5} \draw[reverse arrow=.6, scale = 1] (0,0) to[out=-90,in=-90,looseness=1.5] (1.5,0);
\end{pic}}}}}

\newcommand{\tinycapswap}{\smash{\raisebox{-3pt}{\hspace{-2pt}\ensuremath{\begin{pic}[scale=0.2, yscale=-1]%,string]
   \pgftransformscale{1.5} \draw[reverse arrow=.6, scale = 1] (0,0) to[out=-90,in=-90,looseness=1.5] (1.5,0);
\end{pic}}}}}

\newcommand{\tinyswap}{
\smash{\raisebox{-1pt}{\hspace{-2pt}\ensuremath{
   \begin{pic}[scale=.25]
  \draw (0,-.5) to[out=80,in=-100] (1,.5);
  \draw (1,-.5) to[out=100,in=-80] (0,.5);
\end{pic}}}}
}
%Chris hardcoded picture. Added scale = 2.0 to make up for Aleks' scale = 0.5 
\newenvironment{pic}[1][] {\begin{aligned}\begin{tikzpicture}[scale=2.0, font=\tiny,#1]}{\end{tikzpicture}\end{aligned}} %Hard coded picture 

% \newcommand{\tinycup}{
% \smash{\raisebox{-2pt}{\hspace{-3pt}\ensuremath{\begin{pic}[xscale=0.2,yscale=0.2]
%         \node (10) at (3.75, 0.25) {};
%         \node (12) at (5.25, 0.25) {};
%         \draw [diredge=<, in=-90, out=-90, looseness=2.00] (10.center) to (12.center);
% \end{pic}
% }\hspace{-1pt}}}}

% \newcommand{\tinycap}{
% \smash{\raisebox{-2pt}{\hspace{-3pt}\ensuremath{\begin{pic}[xscale=-0.2,yscale=-0.2]
%         \node (10) at (3.75, 0.25) {};
%         \node (12) at (5.25, 0.25) {};
%         \draw [diredge=<, in=-90, out=-90, looseness=2.00] (10.center) to (12.center);
% \end{pic}
% }\hspace{-1pt}}}}


\newcommand{\tinyring}{
\smash{\raisebox{-2pt}{\hspace{-3pt}\ensuremath{\begin{pic}[scale=0.5]
        \draw[diredge] (1.4,0.5) to[out=-90,in=0] (1.15,0.25) to[out=180,in=-90] (0.9,0.5) to[out=90,in=180] (1.15,0.75) to[out=0,in=90] (1.4,0.5);
\end{pic}
}\hspace{-1pt}}}}

%Small discard symbol. 
\newif\ifvflip\pgfkeys{/tikz/vflip/.is if=vflip}
\newif\ifhflip\pgfkeys{/tikz/hflip/.is if=hflip}
\newif\ifhvflip\pgfkeys{/tikz/hvflip/.is if=hvflip}

\newenvironment{picc}[1][]
{\begin{aligned}\begin{tikzpicture}[font=\tiny,#1]}
{\end{tikzpicture}\end{aligned}}

%\pgfdeclarelayer{foreground}
%\pgfdeclarelayer{background}
%\pgfdeclarelayer{morphismlayer}
%\pgfsetlayers{background,main,morphismlayer,foreground}

\newlength\minimummorphismwidth
\setlength\minimummorphismwidth{0.3cm}
\newlength\stateheight
\setlength\stateheight{0.6cm}
\newlength\minimumstatewidth
\setlength\minimumstatewidth{0.89cm}
\newlength\connectheight
\setlength\connectheight{0.5cm}
\tikzset{colour/.initial=white}

\tikzstyle{pure}=[line width=.7pt]

\makeatletter


\pgfdeclareshape{groundd}
{
    \savedanchor\centerpoint
    {
        \pgf@x=0pt
        \pgf@y=0pt
    }
    \anchor{center}{\centerpoint}
    \anchorborder{\centerpoint}

%     \saveddimen\myscale
%     {
%       \pgfkeysgetvalue{/pgf/scale}{\minwidth}
%       \pgf@x=\minwidth
%     }
    \anchor{north}
    {
        \pgf@x=0pt
        \pgf@y=0.16\stateheight
    }
    \anchor{south}
    {
        \pgf@x=0pt
        \pgf@y=0pt
    }
    \saveddimen\overallwidth
    {
        \pgfkeysgetvalue{/pgf/minimum width}{\minwidth}
        \pgf@x=\minimumstatewidth
        \ifdim\pgf@x<\minwidth
            \pgf@x=\minwidth
        \fi
    }
    \backgroundpath
    {
        \begin{pgfonlayer}{main} %Sean: this used to be 'foreground', swtich back if trouble
        \pgfsetstrokecolor{black}
        \pgfsetlinewidth{1.25pt}
        \ifhflip
            \pgftransformyscale{-1}
        \fi
        \pgftransformscale{0.5}
        \pgfpathmoveto{\pgfpoint{-0.5*\overallwidth}{0pt}}
        \pgfpathlineto{\pgfpoint{0.5*\overallwidth}{0pt}}
        \pgfpathmoveto{\pgfpoint{-0.33*\overallwidth}{0.33*\stateheight}}
        \pgfpathlineto{\pgfpoint{0.33*\overallwidth}{0.33*\stateheight}}
        \pgfpathmoveto{\pgfpoint{-0.16*\overallwidth}{0.66*\stateheight}}
        \pgfpathlineto{\pgfpoint{0.16*\overallwidth}{0.66*\stateheight}}
        \pgfpathmoveto{\pgfpoint{-0.02*\overallwidth}{\stateheight}}
        \pgfpathlineto{\pgfpoint{0.02*\overallwidth}{\stateheight}}
        \pgfusepath{stroke}
        \end{pgfonlayer}
    }
}


%Tiny discard symbol. Was [scale=0.4] after \begin{pic}. Changed to make compatible with TikZit.
\newcommand{\tinyground}[1][groundd]{
\smash{\raisebox{-2pt}{\hspace{-3pt}\ensuremath{\begin{picc}[scale=0.8] 
    \node[#1, scale=0.6] (1) at (0,0.4) {};
    \draw [pure] (1.south) to +(0,-.3);
\end{picc}
}\hspace{-1pt}}}}



\newcommand{\tinygroundupside}[1][groundd]{
\smash{\raisebox{3pt}{\hspace{-3pt}\ensuremath{\begin{picc}[xscale=0.8, yscale=-0.8] 
    \node[#1, yscale=-0.6, xscale=0.6] (1) at (0,0.4) {};
    \draw [pure] (1.south) to +(0,-.3);
\end{picc}
}\hspace{-1pt}}}}

%\newcommand{\discard}[1]{\ensuremath{\tinyground_{#1}}}


%=====================================================
% EVERYTHING BELOW IS FROM OLD FILE PREAMBLE_TIKZIT
%=====================================================

%
% Tikzit Diagrams
%
\usepackage{tikz,xypic}
\usetikzlibrary{decorations.pathreplacing,decorations.markings,arrows.meta,backgrounds,shapes}
\usetikzlibrary{circuits.ee.IEC}
\pgfdeclarelayer{edgelayer}
\pgfdeclarelayer{nodelayer}
\pgfsetlayers{background,edgelayer,nodelayer,main}
\tikzstyle{none}=[inner sep=0mm]
\tikzstyle{every loop}=[]
\tikzstyle{mark coordinate}=[inner sep=0pt,outer sep=0pt,minimum size=3pt,fill=black,circle]

%Warning: this is also defined in the main preamble
\tikzset{arrow/.style={decoration={
    markings,
    mark=at position #1 with \arrow{>[length=2pt, width=3pt]}},
    postaction=decorate},
    reverse arrow/.style={decoration={
    markings,
    mark=at position #1 with {{\arrow{<[length=2pt, width=3pt]}}}},
    postaction=decorate}
}


% GROUND 
\tikzstyle{upground}=[circuit ee IEC,thick,ground,rotate=90,scale=1.5]
\tikzstyle{upgroundwhite}=[circuit ee IEC,thick,ground,rotate=90,scale=1.5, fill=white]
\tikzstyle{downground}=[circuit ee IEC,thick,ground,rotate=-90,scale=1.5]
\tikzstyle{downgroundnorm}=[circuit ee IEC,thick,ground,rotate=-90,scale=1.5, fill=white]
% \tikzstyle{env}=[copoint,regular polygon rotate=0,minimum width=0.2cm, fill=black]
% \tikzstyle{bigground}=[regular polygon,regular polygon sides=3,draw=gray,scale=0.50,inner sep=-0.5pt,minimum width=10mm,fill=gray]
% \tikzstyle{probs}=[shape=semicircle,fill=white,draw=black,shape border rotate=180,minimum width=1.2cm]
%\tikzstyle{grounddag}=[regular polygon,regular polygon sides=3,draw=gray,scale=0.50,inner sep=-0.5pt,minimum width=5mm,fill=gray,regular polygon rotate=180]
%\tikzstyle{ground}=[regular polygon,regular polygon sides=3,draw=gray,scale=0.50,inner sep=-0.5pt,minimum width=5mm,fill=gray]



\newcommand{\mapminh}{5mm} % Aleks and Bob have 6mm
\newcommand{\stateminh}{5mm}
\newcommand{\maplw}{0.7pt} % Alesk and Bob previously just left as default
\newcommand{\stateshift}{-0.2pt}%{1pt}
\newcommand{\effectshift}{-0.2pt}%{1pt}

% BOXES - USED TO BE USED IN CHAPTER 4.
\tikzstyle{box}=[map]
\tikzstyle{medium box}=[medium map]
% \tikzstyle{box}=[draw,shape=rectangle,inner sep=2pt,minimum height=\mapminh,minimum width=6mm,fill=white] %Basic Box with no corners
% \tikzstyle{medium box}=[draw,shape=rectangle,inner sep=2pt,minimum height=\mapminh,minimum width=10mm,fill=white] %Added by Sean
\tikzstyle{dot}=[inner sep=0mm,minimum width=2mm,minimum height=2mm,draw,shape=circle]  
\tikzstyle{black dot}=[dot,fill=black]
\tikzstyle{white dot}=[dot,fill=white,,text depth=-0.2mm]
\tikzstyle{grey dot}=[dot,fill=black!25] %Added by Sean

\tikzstyle{corner1}=[box,fill=white, font=\footnotesize] %
\tikzstyle{corner2}=[dot,fill=white, font=\footnotesize] %
\tikzstyle{corner3}=[dot,fill=black!25, font=\footnotesize] %
\tikzstyle{corner4}=[dot,fill=black, font=\footnotesize] %



%OLD FORM OF SCALARS AS DIAMONDS 
%\tikzstyle{scalar}=[diamond,draw,inner sep=0.5pt,font=\small] %Consider changing to a circle!

\tikzstyle{scalar}=[circle,draw,inner sep=2pt, line width=\maplw] %,font=\small] %Consider changing to a circle!

% \tikzstyle{dscalar}=[diamond,doubled, draw,inner sep=0.5pt,font=\small]

\usetikzlibrary{shapes.misc, positioning}

%Curved boxes for OPT style notation - by Sean
\tikzset{stateshape/.style={append after command={
   \pgfextra
        \draw[sharp corners, fill=white, line width = \maplw]% 
    (\tikzlastnode.west)% 
    [rounded corners=0pt] |- (\tikzlastnode.north)% 
    [rounded corners=0pt] -| (\tikzlastnode.east)% 
    [rounded corners=5pt] |- (\tikzlastnode.south)% 
    [rounded corners=5pt] -| (\tikzlastnode.west);
   \endpgfextra}}}

\tikzset{effectshape/.style={append after command={
   \pgfextra
        \draw[sharp corners, fill=white, line width = \maplw]% 
    (\tikzlastnode.west)% 
    [rounded corners=0pt] |- (\tikzlastnode.south)% 
    [rounded corners=0pt] -| (\tikzlastnode.east)% 
    [rounded corners=5pt] |- (\tikzlastnode.north)% 
    [rounded corners=5pt] -| (\tikzlastnode.west);
   \endpgfextra}}}

%Tweaking state and effect styles. 
 \tikzstyle{map}=[draw,shape=rectangle, inner sep=2pt,minimum height=\mapminh, minimum width=7mm,fill=white]

\tikzstyle{point}=[stateshape,inner sep=2pt, minimum width=6mm, minimum height=\stateminh, yshift=\stateshift]
\tikzstyle{copoint}=[effectshape,inner sep=.2pt, minimum width=6mm, minimum height=\stateminh, yshift=-\effectshift]

\tikzstyle{wide point}=[point, minimum width=12mm]
\tikzstyle{wide copoint}=[copoint, minimum width=12mm]

% fill=white,draw,shape=isosceles triangle,shape border rotate=-90,isosceles triangle stretches=true,inner sep=0pt,minimum width=1.5cm,minimum height=6.12mm,yshift=-0.0mm]


% \tikzstyle{point}=[regular polygon,regular polygon sides=3,draw,scale=0.75,inner sep=-0.5pt,minimum width=9mm,fill=white,regular polygon rotate=180]


% \tikzstyle{small dot}=[inner sep=0.5mm,minimum width=0pt,minimum height=0pt,draw,shape=circle]
% \tikzstyle{small black dot}=[small dot,fill=black]
% \tikzstyle{small white dot}=[small dot,fill=white]


% \tikzstyle{copoint}=[regular polygon,regular polygon sides=3,draw,scale=0.75,inner sep=-0.5pt,minimum width=9mm,fill=white]
% \tikzstyle{dpoint}=[point,doubled]
% \tikzstyle{dcopoint}=[copoint,doubled]

% \tikzstyle{wide copoint}=[fill=white,draw,shape=isosceles triangle,shape border rotate=90,isosceles triangle stretches=true,inner sep=0pt,minimum width=1.5cm,minimum height=6.12mm]



\tikzstyle{decomp}=[fill=white,draw,shape=isosceles triangle,shape border rotate=-90,isosceles triangle stretches=true,inner sep=0pt,minimum width=0.75cm,minimum height=4mm,yshift=-0.0mm]

\tikzstyle{decompwide}=[fill=white,draw,shape=isosceles triangle,shape border rotate=-90,isosceles triangle stretches=true,inner sep=0pt,minimum width=1.5cm,minimum height=4mm,yshift=-0.0mm]

\tikzstyle{decompflip}=[fill=white,draw,shape=isosceles triangle,shape border rotate=90,isosceles triangle stretches=true,inner sep=0pt,minimum width=0.75cm,minimum height=4mm,yshift=-0.0mm]

\tikzstyle{decompwideflip}=[fill=white,draw,shape=isosceles triangle,shape border rotate=90,isosceles triangle stretches=true,inner sep=0pt,minimum width=1.5cm,minimum height=4mm,yshift=-0.0mm]

% \tikzstyle{wide point}=[fill=white,draw,shape=isosceles triangle,shape border rotate=-90,isosceles triangle stretches=true,inner sep=0pt,minimum width=1.5cm,minimum height=6.12mm,yshift=-0.0mm]

% \tikzstyle{wide point plus}=[fill=white,draw,shape=isosceles triangle,shape border rotate=-90,isosceles triangle stretches=true,inner sep=0pt,minimum width=1.74cm,minimum height=7mm,yshift=-0.0mm]

% \tikzstyle{wide dpoint}=[fill=white,doubled,draw,shape=isosceles triangle,shape border rotate=-90,isosceles triangle stretches=true,inner sep=0pt,minimum width=1.5cm,minimum height=6.12mm,yshift=-0.0mm]

% \tikzstyle{tinypoint}=[regular polygon,regular polygon sides=3,draw,scale=0.55,inner sep=-0.15pt,minimum width=6mm,fill=white,regular polygon rotate=180] 

 \tikzstyle{map}=[draw,shape=rectangle, inner sep=2pt,minimum height=\mapminh, minimum width=7mm,fill=white, line width = \maplw]

\tikzstyle{medium map} = [map, minimum width = 12mm] 
\tikzstyle{semilarge map} = [map, minimum width = 15mm] 
\tikzstyle{large map} = [map, minimum width = 18mm] 

% Use non-dagger states as standard for the following two.

\tikzstyle{kpoint} =[point]
\tikzstyle{kpointadj} =[copoint]
\tikzstyle{kpointconj}=[dagpointconj] %Need daggers for conjugates
% \tikzstyle{dagpointtrans}=[shape=cornercopoint,shorten right=5pt,dagpoint common]
% \tikzstyle{dagpointsymm}=[shape=cornerpoint,shorten left=5pt,shorten right=5pt,dagpoint common]



% \tikzstyle{state}=[point]%stateshape,inner sep=.4ex, node on layer=foreground, minimum width = 0.6cm, minimum height = 0.4cm, font = \small]

% \tikzstyle{effect}=[copoint]




% ==== ALEKS AND BOB POINTY BOXES === call this dagmap

\makeatletter
\newcommand{\boxshape}[3]{%
\pgfdeclareshape{#1}{
\inheritsavedanchors[from=rectangle] % this is nearly a rectangle
\inheritanchorborder[from=rectangle]
\inheritanchor[from=rectangle]{center}
\inheritanchor[from=rectangle]{north}
\inheritanchor[from=rectangle]{south}
\inheritanchor[from=rectangle]{west}
\inheritanchor[from=rectangle]{east}
% ... and possibly more
\backgroundpath{% this is new
% store lower right in xa/ya and upper right in xb/yb
\southwest \pgf@xa=\pgf@x \pgf@ya=\pgf@y
\northeast \pgf@xb=\pgf@x \pgf@yb=\pgf@y

\@tempdima=#2
\@tempdimb=#3

\pgfpathmoveto{\pgfpoint{\pgf@xa - 5pt + \@tempdima}{\pgf@ya}}
\pgfpathlineto{\pgfpoint{\pgf@xa - 5pt - \@tempdima}{\pgf@yb}}
\pgfpathlineto{\pgfpoint{\pgf@xb + 5pt + \@tempdimb}{\pgf@yb}}
\pgfpathlineto{\pgfpoint{\pgf@xb + 5pt - \@tempdimb}{\pgf@ya}}
\pgfpathlineto{\pgfpoint{\pgf@xa - 5pt + \@tempdima}{\pgf@ya}}
\pgfpathclose
}
}}

%Aleks had the boxes at 5pt differences before, e.g. NEbox was {0pt}{5pt}
% \boxshape{NEbox}{0pt}{5pt}
% \boxshape{SEbox}{0pt}{-5pt}
% \boxshape{NWbox}{5pt}{0pt}
% \boxshape{SWbox}{-5pt}{0pt}
\boxshape{NEbox}{0pt}{3pt} %My setting of 3pt
\boxshape{SEbox}{0pt}{-3pt}
\boxshape{NWbox}{3pt}{0pt}
\boxshape{SWbox}{-3pt}{0pt}
\boxshape{rec-box}{0pt}{0pt}
%\boxshape{EBox}{-3pt}{3pt}
%\boxshape{WBox}{3pt}{-3pt}
\makeatother

\tikzstyle{cloud}=[shape=cloud,draw,minimum width=1.5cm,minimum height=1.5cm]




% [draw,shape=rectangle,inner sep=2pt,minimum height=\mapminh,minimum width=6mm,fill=white]


%draw,shape=NEbox,inner sep=2pt,minimum height=\mapminh,fill=white,minimum width=7mm, line width = \maplw]

 % box, draw,line width = \maplw, inner sep=2pt, minimum height=5mm, minimum width=7mm] %

% \tikzstyle{map}=[draw,shape=rec-box,inner sep=2pt,minimum height=\mapminh,
% minimum width=3mm, %Sean: I added this, copied it from 'box'
% fill=white, line width = \maplw] %

% \tikzstyle{box}=[draw,shape=rectangle,inner sep=2pt,minimum height=\mapminh,minimum width=6mm,fill=white] %Basic Box with no corners





\tikzstyle{dagmap}=[draw,shape=NEbox,inner sep=2pt,minimum height=\mapminh,fill=white, line width = \maplw] %
\tikzstyle{dashedmap}=[draw,dashed,shape=NEbox,inner sep=2pt,minimum height=\mapminh,fill=white, line width = \maplw]
\tikzstyle{mapdag}=[draw,shape=SEbox,inner sep=2pt,minimum height=\mapminh,fill=white, line width = \maplw]
\tikzstyle{mapadj}=[draw,shape=SEbox,inner sep=2pt,minimum height=\mapminh,fill=white, line width = \maplw]
\tikzstyle{maptrans}=[draw,shape=SWbox,inner sep=2pt,minimum height=\mapminh,fill=white, line width = \maplw]
\tikzstyle{mapconj}=[draw,shape=NWbox,inner sep=2pt,minimum height=\mapminh,fill=white, line width = \maplw]


\tikzstyle{medium dagmap}=[draw,shape=NEbox,inner sep=2pt,minimum height=\mapminh,fill=white,minimum width=7mm, line width = \maplw]
% \tikzstyle{medium map dag}=[draw,shape=SEbox,inner sep=2pt,minimum height=\mapminh,fill=white,minimum width=7mm, line width = \maplw]
% \tikzstyle{medium map adj}=[draw,shape=SEbox,inner sep=2pt,minimum height=\mapminh,fill=white,minimum width=7mm, line width = \maplw]
% \tikzstyle{medium map trans}=[draw,shape=SWbox,inner sep=2pt,minimum height=\mapminh,fill=white,minimum width=7mm, line width = \maplw]
% \tikzstyle{medium map conj}=[draw,shape=NWbox,inner sep=2pt,minimum height=\mapminh,fill=white,minimum width=7mm, line width = \maplw]
\tikzstyle{semilarge dagmap}=[draw,shape=NEbox,inner sep=2pt,minimum height=\mapminh,fill=white,minimum width=9.5mm, line width = \maplw]
% \tikzstyle{semilarge map trans}=[draw,shape=SWbox,inner sep=2pt,minimum height=\mapminh,fill=white,minimum width=9.5mm, line width = \maplw]
% \tikzstyle{semilarge map adj}=[draw,shape=SEbox,inner sep=2pt,minimum height=\mapminh,fill=white,minimum width=9.5mm, line width = \maplw]
% \tikzstyle{semilarge map dag}=[draw,shape=SEbox,inner sep=2pt,minimum height=\mapminh,fill=white,minimum width=9.5mm, line width = \maplw]
% \tikzstyle{semilarge map conj}=[draw,shape=NWbox,inner sep=2pt,minimum height=\mapminh,fill=white,minimum width=9.5mm, line width = \maplw]
\tikzstyle{large dagmap}=[draw,shape=NEbox,inner sep=2pt,minimum height=\mapminh,fill=white,minimum width=12mm, line width = \maplw]
% \tikzstyle{large map conj}=[draw,shape=NWbox,inner sep=2pt,minimum height=\mapminh,fill=white,minimum width=12mm, line width = \maplw]
% \tikzstyle{very large map}=[draw,shape=NEbox,inner sep=2pt,minimum height=\mapminh,fill=white,minimum width=17mm, line width = \maplw]



%=========== ALEKS AND BOB STATE SHAPE ('kpoint') ============
\makeatletter

\pgfdeclareshape{cornerpoint}{
\inheritsavedanchors[from=rectangle] % this is nearly a rectangle
\inheritanchorborder[from=rectangle]
\inheritanchor[from=rectangle]{center}
\inheritanchor[from=rectangle]{north}
\inheritanchor[from=rectangle]{south}
\inheritanchor[from=rectangle]{west}
\inheritanchor[from=rectangle]{east}
% ... and possibly more
\backgroundpath{% this is new
% store lower right in xa/ya and upper right in xb/yb
\southwest \pgf@xa=\pgf@x \pgf@ya=\pgf@y
\northeast \pgf@xb=\pgf@x \pgf@yb=\pgf@y

\pgfmathsetmacro{\pgf@shorten@left}{\pgfkeysvalueof{/tikz/shorten left}}
\pgfmathsetmacro{\pgf@shorten@right}{\pgfkeysvalueof{/tikz/shorten right}}

\pgfpathmoveto{\pgfpoint{0.5 * (\pgf@xa + \pgf@xb)}{\pgf@ya - 5pt}}
\pgfpathlineto{\pgfpoint{\pgf@xa - 8pt + \pgf@shorten@left}{\pgf@yb - 1.5 * \pgf@shorten@left}}
\pgfpathlineto{\pgfpoint{\pgf@xa - 8pt + \pgf@shorten@left}{\pgf@yb}}
\pgfpathlineto{\pgfpoint{\pgf@xb + 8pt - \pgf@shorten@right}{\pgf@yb}}
\pgfpathlineto{\pgfpoint{\pgf@xb + 8pt - \pgf@shorten@right}{\pgf@yb - 1.5 * \pgf@shorten@right}}
\pgfpathclose
}
}

\pgfdeclareshape{cornercopoint}{
\inheritsavedanchors[from=rectangle] % this is nearly a rectangle
\inheritanchorborder[from=rectangle]
\inheritanchor[from=rectangle]{center}
\inheritanchor[from=rectangle]{north}
\inheritanchor[from=rectangle]{south}
\inheritanchor[from=rectangle]{west}
\inheritanchor[from=rectangle]{east}
% ... and possibly more
\backgroundpath{% this is new
% store lower right in xa/ya and upper right in xb/yb
\southwest \pgf@xa=\pgf@x \pgf@ya=\pgf@y
\northeast \pgf@xb=\pgf@x \pgf@yb=\pgf@y

\pgfmathsetmacro{\pgf@shorten@left}{\pgfkeysvalueof{/tikz/shorten left}}
\pgfmathsetmacro{\pgf@shorten@right}{\pgfkeysvalueof{/tikz/shorten right}}

\pgfpathmoveto{\pgfpoint{0.5 * (\pgf@xa + \pgf@xb)}{\pgf@yb + 5pt}}
\pgfpathlineto{\pgfpoint{\pgf@xa - 8pt + \pgf@shorten@left}{\pgf@ya + 1.5 * \pgf@shorten@left}}
\pgfpathlineto{\pgfpoint{\pgf@xa - 8pt + \pgf@shorten@left}{\pgf@ya}}
\pgfpathlineto{\pgfpoint{\pgf@xb + 8pt - \pgf@shorten@right}{\pgf@ya}}
\pgfpathlineto{\pgfpoint{\pgf@xb + 8pt - \pgf@shorten@right}{\pgf@ya + 1.5 * \pgf@shorten@right}}
\pgfpathclose
}
}

\makeatother

\pgfkeyssetvalue{/tikz/shorten left}{0pt}
\pgfkeyssetvalue{/tikz/shorten right}{0pt}

%SEAN: THIS USED TO BE CALLED 'kpoint'. Its now called 'dagpoint'

%SEAN: If you add min width = 4mm and min height = 5mm to kpoint common it will make line up with maps. 
\tikzstyle{dagpoint common}=[draw,fill=white,inner sep=1pt, line width = \maplw, minimum height = 4mm, yshift=1.2pt] %SEAN: I ADDED THIS yshift to make it line up with maps. 1pt seems to work.
\tikzstyle{dagpoint sc}=[shape=cornerpoint,dagpoint common]
\tikzstyle{dagpoint adjoint sc}=[shape=cornercopoint,dagpoint common]
\tikzstyle{dagpoint}=[shape=cornerpoint,shorten left=4pt,dagpoint common]
\tikzstyle{dagpointadj}=[shape=cornercopoint,shorten left=5pt,dagpoint common]
\tikzstyle{dagpointconj}=[shape=cornerpoint,shorten right=5pt,dagpoint common]
\tikzstyle{dagpointtrans}=[shape=cornercopoint,shorten right=5pt,dagpoint common]
\tikzstyle{dagpointsymm}=[shape=cornerpoint,shorten left=5pt,shorten right=5pt,dagpoint common]

\tikzstyle{widedagpoint}=[dagpoint, minimum width=1 cm, inner sep=2pt]%, text depth=-0.7 mm]
\tikzstyle{widedagpointadj}=[dagpointadj, minimum width=1 cm, inner sep=2pt]%, text depth=0.7 mm]
% \tikzstyle{dagpointdag}=[dagpoint adjoint]

% \tikzstyle{big kpoint}=[kpoint, minimum width=1.2 cm, minimum height=8mm, inner sep=4pt, text depth=3mm]

% \tikzstyle{wide kpoint}=[kpoint, minimum width=1 cm, inner sep=2pt]%, text depth=-0.7 mm]
% \tikzstyle{wide kpointdag}=[kpointdag, minimum width=1 cm, inner sep=2pt]%, text depth=0.7 mm]
% \tikzstyle{wide kpointconj}=[kpointconj, minimum width=1 cm, inner sep=2pt]%, text depth=-0.7 mm]
% \tikzstyle{wide kpointtrans}=[kpointtrans, minimum width=1 cm, inner sep=2pt]%, text depth=0.7 mm]

% \tikzstyle{gray kpoint}=[kpoint,fill=gray!50!white]
% \tikzstyle{gray kpointdag}=[kpointdag,fill=gray!50!white]
% \tikzstyle{gray kpointadj}=[kpointadj,fill=gray!50!white]
% \tikzstyle{gray kpointconj}=[kpointconj,fill=gray!50!white]
% \tikzstyle{gray kpointtrans}=[kpointtrans,fill=gray!50!white]

% \tikzstyle{gray dkpoint}=[kpoint,fill=gray!50!white,doubled]
% \tikzstyle{gray dkpointdag}=[kpointdag,fill=gray!50!white,doubled]
% \tikzstyle{gray dkpointadj}=[kpointadj,fill=gray!50!white,doubled]
% \tikzstyle{gray dkpointconj}=[kpointconj,fill=gray!50!white,doubled]
% \tikzstyle{gray dkpointtrans}=[kpointtrans,fill=gray!50!white,doubled]

% \tikzstyle{white label}=[draw,fill=white,rectangle,inner sep=0.7 mm]
% \tikzstyle{gray label}=[draw,fill=gray!50!white,rectangle,inner sep=0.7 mm]
% \tikzstyle{black label}=[draw,fill=black,rectangle,inner sep=0.7 mm]

% \tikzstyle{dkpoint}=[kpoint,doubled]
% \tikzstyle{wide dkpoint}=[wide kpoint,doubled]
% \tikzstyle{dkpointdag}=[kpoint adjoint,doubled]
% \tikzstyle{wide dkpointdag}=[wide kpointdag,doubled]
% \tikzstyle{dkcopoint}=[kpoint adjoint,doubled]
% \tikzstyle{dkpointadj}=[kpoint adjoint,doubled]
% \tikzstyle{dkpointconj}=[kpoint conjugate,doubled]
% \tikzstyle{dkpointtrans}=[kpoint transpose,doubled]

% \tikzstyle{kscalar}=[kpoint common, shape=EBox, inner xsep=-1pt, inner ysep=3pt,font=\small]
% \tikzstyle{kscalarconj}=[kpoint common, shape=WBox, inner xsep=-1pt, inner ysep=3pt,font=\small]

% \tikzstyle{spekpoint}=[kpoint sc,minimum height=5mm,inner sep=3pt]
% \tikzstyle{spekcopoint}=[kpoint adjoint sc,minimum height=5mm,inner sep=3pt]

% \tikzstyle{dspekpoint}=[spekpoint,doubled]
% \tikzstyle{dspekcopoint}=[spekcopoint,doubled]

% %Morphism shape for dagger categories 
% \newif\ifvflip\pgfkeys{/tikz/vflip/.is if=vflip}
% \newif\ifhflip\pgfkeys{/tikz/hflip/.is if=hflip}
% \newif\ifhvflip\pgfkeys{/tikz/hvflip/.is if=hvflip}
% \newlength\morphismheight
% \setlength\morphismheight{4mm}
% \newlength\wedgewidth
% \setlength\wedgewidth{8pt}
% \tikzset{width/.initial=1mm}
% \makeatletter


% \tikzstyle{morphism}=[font=\small,morphismshape] 
% \tikzstyle{morphismflip}=[morphism, hflip]

% \pgfdeclareshape{morphismshape}
% {
%     \savedanchor\centerpoint
%     {
%         \pgf@x=0pt
%         \pgf@y=0pt
%     }
%     \anchor{center}{\centerpoint}
%     \anchorborder{\centerpoint}
%     \saveddimen\overallwidth
%     {
%         \pgfkeysgetvalue{/tikz/width}{\minwidth}
%         \pgf@x=\wd\pgfnodeparttextbox
%         \ifdim\pgf@x<\minwidth
%             \pgf@x=\minwidth
%         \fi
%     }
%     \savedanchor{\upperrightcorner}
%     {
%         \pgf@y=.5\ht\pgfnodeparttextbox
%         \advance\pgf@y by -.5\dp\pgfnodeparttextbox
%         \pgf@x=.5\wd\pgfnodeparttextbox
%     }
%     \anchor{north}
%     {
%         \pgf@x=0pt
%         \pgf@y=0.5\morphismheight
%     }
%     \anchor{north east}
%     {
%         \pgf@x=\overallwidth
%         \multiply \pgf@x by 2
%         \divide \pgf@x by 5
%         \pgf@y=0.5\morphismheight
%     }
%     \anchor{east}
%     {
%         \pgf@x=\overallwidth
%         \divide \pgf@x by 2
%         \advance \pgf@x by 5pt
%         \pgf@y=0pt
%     }
%     \anchor{west}
%     {
%         \pgf@x=-\overallwidth
%         \divide \pgf@x by 2
%         \advance \pgf@x by -5pt
%         \pgf@y=0pt
%     }
%     \anchor{north west}
%     {
%         \pgf@x=-\overallwidth
%         \multiply \pgf@x by 2
%         \divide \pgf@x by 5
%         \pgf@y=0.5\morphismheight
%     }
%     \anchor{south east}
%     {
%         \pgf@x=\overallwidth
%         \multiply \pgf@x by 2
%         \divide \pgf@x by 5
%         \pgf@y=-0.5\morphismheight
%     }
%     \anchor{south west}
%     {
%         \pgf@x=-\overallwidth
%         \multiply \pgf@x by 2
%         \divide \pgf@x by 5
%         \pgf@y=-0.5\morphismheight
%     }
%     \anchor{south}
%     {
%         \pgf@x=0pt
%         \pgf@y=-0.5\morphismheight
%     }
%     \anchor{text}
%     {
%         \upperrightcorner
%         \pgf@x=-\pgf@x
%         \pgf@y=-\pgf@y
%     }
%     \backgroundpath
%     {
%     \begin{scope}
%         \pgfkeysgetvalue{/tikz/fill}{\morphismfill}
%         \pgfsetstrokecolor{black}
%         \pgfsetlinewidth{.7pt}
%         \begin{scope}
%         \pgfsetstrokecolor{black}
%         \pgfsetfillcolor{white}
%         \pgfsetlinewidth{.7pt}
%                 \ifhflip
%                     \pgftransformyscale{-1}
%                 \fi
%                 \ifvflip
%                     \pgftransformxscale{-1}
%                 \fi
%                 \ifhvflip
%                     \pgftransformxscale{-1}
%                     \pgftransformyscale{-1}
%                 \fi
%                 \pgfpathmoveto{\pgfpoint
%                     {-0.5*\overallwidth-5pt}
%                     {0.5*\morphismheight}}
%                 \pgfpathlineto{\pgfpoint
%                     {0.5*\overallwidth+5pt}
%                     {0.5*\morphismheight}}
%                 \pgfpathlineto{\pgfpoint
%                     {0.5*\overallwidth + \wedgewidth}
%                     {-0.5*\morphismheight}}
%                 \pgfpathlineto{\pgfpoint
%                     {-0.5*\overallwidth-5pt}
%                     {-0.5*\morphismheight}}
%                 \pgfpathclose
%                 \pgfusepath{fill,stroke}
%         \end{scope}
%     \end{scope}
%     }
% }
% \makeatother


\tikzstyle{every picture}=[baseline=-0.25em,scale=0.5]
\tikzstyle{label}=[font=\footnotesize,text height=1ex, text depth=0.15ex]



% \tikzstyle{statewide}=[state, minimum width = 1.2cm]
% \tikzstyle{stateextrawide}=[state, minimum width = 1.4cm]


% MULTICATEGORIES

\usetikzlibrary[shapes]


%\usetikzlibrary{shapes.geometric}
%\usepgflibrary{shapes.geometric}
%======================================================
% Put these after other definitions in the Preamble
\tikzset{
sidetriangle/.style = {regular polygon, regular polygon sides = 3, aspect = 1, shape border rotate = 90, draw, inner sep = 0, minimum width = 1.2cm}
}


\tikzset{
isoc/.style = {shape=isosceles triangle, shape border rotate = 180, isosceles triangle stretches = true, minimum width = 1.2cm, minimum height= 1.5cm, inner sep = 0.3}}

%Coarse-graining Style (Prob don't need these).
\tikzset{
coarse/.style = {shape = circle, fill = white, draw, inner sep = 0, minimum width =0.125cm}
}
\tikzset{
coarsesymbol/.style = {shape = circle, fill = white, inner sep = -0.7, minimum width = 0.125cm}
}
%=====================================================



\tikzstyle{sidetriangle2}=[sidetriangle, minimum width = 2cm, fill=white]
\tikzstyle{sideisocsmall}]=[style=isoc, minimum width = 1cm, minimum height = 0.8cm, draw, fill=white, font=\Large]
\tikzstyle{sideisoc}]=[style=isoc, minimum width = 2cm, draw, fill=white, font=\Large]
\tikzstyle{sideisocmid}]=[style=isoc, minimum width = 2.5cm, draw, fill=white, font=\Large]
\tikzstyle{sideisocmedium}]=[style=isoc, minimum width = 3cm, draw, fill=white, font=\Large]


\newcommand{\tinygroundnew}{
\smash{
% \raisebox{-2pt}
{\hspace{-3pt}
\ensuremath{
\begin{picc}[scale=1.0] 
    \node[upground, xscale=0.8, yscale=0.7] (1) at (0,0.16) {};
    \draw (0,0.03) to (0,-0.25);
\end{picc}
}\hspace{-1pt}}}}



\newcommand{\tinygroundflipnew}{
\smash{
% \raisebox{1pt}
{\hspace{-3pt}
\ensuremath{
\begin{picc}[yscale=-1.0] 
    \node[upground, xscale=0.8, yscale=-0.7] (1) at (0,0.10) {};
    \draw (0,-0.03) to (0,-0.31);
\end{picc}
}\hspace{-1pt}}}}


\newcommand{\tinygroundflipnewer}{
\smash{
% \raisebox{1pt}
{\hspace{-3pt}
\ensuremath{
\begin{picc}[yscale=-1.0] 
    \node[upgroundwhite, xscale=0.8, yscale=-0.7] (1) at (0,0.10) {};
    \draw (0,-0.03) to (0,-0.31);
\end{picc}
}\hspace{-1pt}}}}

\newcommand{\uniformprob}{
\smash{
% \raisebox{1pt}
{\hspace{-4pt}
\ensuremath{
\begin{aligned}\begin{tikzpicture}[]
%
    \node[scale=1] at (0,-0.15) {p};
	\draw[line width=0.2mm] (-.13,0.15) to (0.09,0.15);
	\draw[line width=0.2mm] (-.09,0.22) to (0.05,0.22);
	\draw[line width=0.2mm] (-.06,0.29) to (0.02,0.29);
%  
%	\draw[line width=0.2mm] (-.09,-.05) to (0.13,-.05);
%	\draw[line width=0.2mm] (-.09,-.115) to (0.13,-.115);
%
%	\draw[line width=0.2mm] (-.09,-.05) to (0.13,-.05);
%	\draw[line width=0.2mm] (-.05,-.115) to (0.09,-.115);
%	\draw[line width=0.2mm] (-.025,-.18) to (0.055,-.18);
%
%	\draw (-.09,-.05) to (0.13,-.05);
%	\draw (-.05,-.1) to (0.09,-.1);
%	\draw (-.09,-.15) to (0.13,-.15);	
%
	%	\node[circuit ee IEC,thick,ground,rotate=90,scale=.8, xscale=1, yscale=-0.8] (1) at (-.02,0.23) {};
	%\node[circuit ee IEC,thick,ground,rotate=90,scale=.6, xscale=1, yscale=-0.8] (1) at (-.01,-.1) {};
\end{tikzpicture}\end{aligned}
}
\hspace{-5pt}}}}

\newcommand{\Cond}[1]{{\scriptstyle |}#1{\scriptstyle \ra}}
\newcommand{\Mar}[1]{{\scriptstyle \la #1 |}}
\newcommand{\Bot}{{\scriptscriptstyle \bot}}
\newcommand{\Botl}{{\!\scriptscriptstyle \bot}}
%\tikzset{stateshaperot/.style={append after command={
%   \pgfextra
%        \draw[sharp corners, fill=white, line width = 0.15mm]% 
%    (\tikzlastnode.west)% 
%    [rounded corners=0pt] |- (\tikzlastnode.north)% 
%    [rounded corners=2pt] -| (\tikzlastnode.east)% 
%    [rounded corners=2pt] |- (\tikzlastnode.south)% 
%    [rounded corners=0pt] -| (\tikzlastnode.west);
%   \endpgfextra}}}
%\tikzstyle{pointrot}=[stateshaperot,inner sep=1pt, minimum width=0mm, minimum height=0mm, yshift=0mm]
%\newcommand{\Cond}[1]{
%\begin{tikzpicture}
%	\begin{pgfonlayer}{nodelayer}
%		\node [scale=1,style=pointrot] at (0, -0.0) {$#1$};
%	\end{pgfonlayer}
%\end{tikzpicture}
%}
